% Draft section: concurrency demonstration (schedules vs traces).
%
% Standalone fragment for the POPL paper. Citation keys are placeholders;
% \cref targets refer to sections of the eventual paper. Every number is
% measured: reproduce with `cargo run -p alifib --release --example
% trace_search` against examples/{Philosophers,Corridor,PhilosophersRing}.ali
% (commits 56cd92c, 9b1ffa6).

\section{Traces, not schedules}
\label{section:traces}

The purpose of this section is to put the trace objects of
\cref{section:transparency} to work on a problem they were not designed
for: the analysis of concurrent systems. alifib has no concurrency
features---no processes, no channels, no independence relation, and no
notion of commutation. I will argue that it needs none of them. The
structure that concurrency theory works hardest to recover is already
present in the geometry of the engine's values, and the standard
apparatus---trace equivalence, partial-order reduction, deadlock
detection---falls out of its ordinary notion of equality.

\paragraph{Schedules and traces.}
Consider the simplest possible concurrent system: a word $a\,b$, a rule
$r_a : a \to x$, and a rule $r_b : b \to y$. The two redexes are
disjoint; neither needs the other. A \emph{schedule} is a run as it is
actually performed: a sequence of choices, $r_a$ then $r_b$, or $r_b$
then $r_a$. This is what nearly every model of computation calls an
execution, and the order is part of the data even when it is
meaningless. A \emph{trace}\footnote{The word is overloaded. Mazurkiewicz
traces are equivalence classes of sequences; the traces of this paper
are diagrams, and their relation to the former is precisely the point of
this section.} is the diagram that the run assembles: a $2$-cell
labelled $r_a$ and a $2$-cell labelled $r_b$, sitting side by side,
because that is where the redexes are. There is no slot in this object
that could record which cell was applied first. Both schedules build the
same diagram, and the engine agrees: started on one and aimed at the
other, it reports the target reached before a single step is taken. In
general, a schedule is a total order on the events of a run; the trace
records exactly the causal one.

\paragraph{Mazurkiewicz equivalence, definitionally.}
The classical remedy for the spurious total order is due to
\textcite{mazurkiewicz_1977}: fix an independence relation on actions,
and quotient sequences by the congruence that commutes independent
neighbours. The remedy presupposes the disease. One must first commit to
sequences, then specify---as additional data, subject to side
conditions---which reorderings are harmless. Here the situation is
inverted. The trace is the only object the engine ever builds; the
interleavings are merely its \emph{layerings}, recoverable when wanted
and never primitive. The slogan: \emph{the trace does not forget the
interleaving; it never recorded it.} Independence, too, is not declared
but read off: two available rewrites commute just when their matches are
disjoint subdiagrams, and where they overlap---two philosophers reaching
for one fork---the traces genuinely branch. What an event structure
posits as axioms, the diagram exhibits as geometry.

\paragraph{Two benchmarks.}
I test this on the two standard topologies of the folklore. The first is
\emph{dining philosophers in a row}: forks and philosophers alternate
along a word, every philosopher chooses which adjacent fork to take
first, and neighbours race for the fork they share. Five rules suffice,
and---a pleasing consequence of data being words---the type does not
mention the number of philosophers at all: the instance is the initial
word $f\,p\,f\,p\,\cdots\,f$. The second is a \emph{corridor}: a single
track of $m$ free segments claimed from opposite ends by two trains, or
equivalently two transactions acquiring a chain of locks in opposite
orders. The corridor is doomed by design; whoever wins the race for the
last segment, the next state is a head-on meeting at which no rule
applies. Its combinatorics are exact: $2^m$ schedules (each segment goes
to one train or the other), but only $m+1$ traces---one causal story per
meeting point.

\paragraph{The search.}
A driver of some two hundred lines explores these systems through the
same session interface a human uses. It enumerates the reachable states
by breadth-first search, deduplicating by the engine's isomorphism;
counts schedules exactly, by path-counting over the resulting acyclic
graph; and classifies runs into trace classes by walking every schedule
with step and undo, collapsing the assembled proof diagrams under that
same isomorphism. Notice what the driver does \emph{not} contain: no
independence relation, no commutation check, no persistent sets. The
entire quotient is one call to the equality of values. The measurements:

\begin{center}
\begin{tabular}{l r r r l}
\toprule
system & states & schedules (exact) & traces & verdict \\
\midrule
philosophers, $N=4$ & $377$ & $119{,}328$ & $128$ & deadlock-free \\
philosophers, $N=5$ & $1{,}597$ & $30{,}700{,}608$ & $\geq 220$ & deadlock-free \\
corridor, $m=10$ & $66$ & $1{,}024$ & $11$ & $11$ deadlocks \\
corridor, $m=40$ & $861$ & $2^{40} \approx 1.1 \times 10^{12}$ & --- & $41$ deadlocks \\
round table, $N=3$ & $1{,}574$ & $\infty$ (cyclic) & --- & $2$ deadlocks \\
\bottomrule
\end{tabular}
\end{center}

At $N=4$ every one of the $119{,}328$ schedules is walked, and they
collapse to $128$ certificates: a quotient of roughly $930$, measured
rather than asserted. As a sanity check, the corridor at $m=10$ yields
eleven classes whose schedule counts are $252, 210, 210, 120, 120, 45,
45, 10, 10, 1, 1$---precisely the tenth row of Pascal's triangle. The
brute-force walk recomputes the closed form, which is some evidence that
it is walking what we think it is walking. Each deadlock arrives as a
certificate, an ordinary diagram one can replay step by step; at $m=4$,
in full:
\[
  (w\, s\, s\, s\, \mathit{adv}_e)
  \mathbin{\#_1} (w\, s\, s\, \mathit{adv}_e\, s)
  \mathbin{\#_1} (\mathit{adv}_w\, s\, e\, s\, s)
  \mathbin{\#_1} (s\, \mathit{adv}_w\, e\, s\, s).
\]
The two trains advance, meet, and stop; the diagram \emph{is} the
space--time picture of the collision.

\paragraph{Exhaustive verdicts.}
Walking schedules is exponential, and past $150{,}000$ of them the
driver merely samples---whence the lower bound at $N=5$. But the
verification claims never rest on sampling. The reachable state graph is
finite whatever the schedule count, and a terminal state of the graph is
terminal for every run that reaches it; so the verdict column of the
table is exhaustive even at $2^{40}$, where all $41$ meeting points are
found without any pretence of walking a trillion schedules. Safety over
all interleavings is a statement about a small graph.

\paragraph{The round table.}
The reader will have noticed that the philosophers above sit in a row,
and the classical deadlock needs a circle. The obstruction is real: the
values of alifib in dimension one are molecules, molecules are
intervals, and a cyclic word is simply not a value. The circle cannot
live in the geometry of the word. It can, however, live in the
\emph{typing}. Give each fork its own sort, fix the philosophers in
place, and let free forks drift past everything by structural mobility
rules; then fork $f_i$ is shared by philosophers $i$ and $i-1$ because
only their take-rules mention it, and the wait-for cycle returns at
$N=3$ with indices read modulo $3$. The search finds what it should:
$1{,}574$ reachable states; exactly two terminal ones, all-left and
all-right---the two circular waits, discovered, not planted; and the
basin from which success is unreachable is exactly those two states, so
the folklore theorem---the round table deadlocks just when every
philosopher grabs in the same direction, and recovery is possible from
anywhere else---is verified mechanically. The encoding is not free of
charge. The drift rules are \emph{structural}: they make the state
graph cyclic, the schedule space infinite, and the trace count a
question about rewriting \emph{modulo} the structural cells---the same
obligation that the term-rewriting encoding of \cref{section:trs} faces
for its own structural layer, and open for the same reason. I prefer to
state the frontier than to hide it: where the free theory ends is
exactly where the modulo-theory must begin.

\paragraph{What this does and does not show.}
At this point a reader from model checking may protest that these are
classroom benchmarks; that SPIN visits more states per second than this
engine will visit in its lifetime \parencite{holzmann_2003}; and that
partial-order reduction is thirty years of hard-won
engineering \parencite{godefroid_1996, flanagan_godefroid_2005}, not a
table with five rows. So much is true. But the claim here is not speed,
and it is not scale. The claim is that the object those algorithms
labour to construct---a canonical representative of each commutation
class---is, in this setting, the only object that exists, and that
``two runs are equivalent'' is decided by the same isomorphism that
decides every other equality of values. Persistent sets and sleep sets
are answers to a question that this syntax does not ask.
\textcite{pratt_1991} and \textcite{vanglabbeek_2006} argued long ago
that concurrency is inherently higher-dimensional; what their
higher-dimensional automata lacked was not insight but an executable
syntax with a matching engine. And on the Petri-net side the
correspondence is older still: the trace diagram of a run is its causal
process in the sense of \textcite{goltz_reisig_1983}, computed here by
running the net rather than defined over it---one wire-swap short, as
above, of the symmetric monoidal semantics of
\textcite{meseguer_montanari_1990}. The sum of it is this: the partial
order was never the quotient of the total ones. It was the datum all
along.
